% f2000_spring_jingwan_science_q20_fig2.tex — 2000 年·春季京皖卷·理科 第 20 题 图二（折叠后三棱锥 D′-ABC，直二面角）
%
% 【来源】`.work/phase4-b20/source/2000春季京皖理.pdf`
%   （SHA-256 37f07244497aa0a21be12d8a31eac78b61490f835f088719a5990463585a35ce），
%   第 1 页（印刷页 61）；qpdf --qdf 解包后对内容流逐路径解析。
% 【坐标口径】源 PDF 用户空间（原点左下、y 向上）。本图本地坐标
%   (x, y) = (x_page − 949.762, y_page − 629.419)，1 单位 = 1 TeX pt；
%   A 取底面左端点；AB 水平（y_page ≡ 629.419）。
% 【元素计数】6 条路径（AB／BD′／D′A 实线；AC／CB／CD′ 虚线）＋ 4 个标签。
% 【标签】位置取 pdftotext -bbox 文本中心换算（源中心逐条记在 \node 行末）；
%   本图有 1 处微调：「C」下移 3.0 pt，避开 AC 虚线（原净距 2.0 pt）；
%   其余标签按源中心放置。
\begin{tikzpicture}[x=1pt, y=1pt,
    every node/.style={inner sep=0pt, font=\fontsize{10.5pt}{12.6pt}\selectfont}]
  \draw[line width=0.399pt] (0,0) -- (68.037,0);
  \draw[line width=0.399pt] (68.037,0) -- (17.009,34.019);
  \draw[line width=0.399pt] (17.009,34.019) -- (0,0);
  \draw[line width=0.399pt, dash pattern=on 2.99pt off 2.99pt] (0,0) -- (20.411,20.411);
  \draw[line width=0.399pt, dash pattern=on 2.99pt off 2.99pt] (68.037,0) -- (20.411,20.411);
  \draw[line width=0.399pt, dash pattern=on 2.99pt off 2.99pt] (20.411,20.411) -- (17.009,34.019);

  \node[anchor=center] at (-4.932,-5.889) {$A$};  % 源中心 (944.83,623.53)
  \node[anchor=center] at (72.708,-5.889) {$B$};  % 源中心 (1022.47,623.53)
  \node[anchor=center] at (20.038,9.471) {$C$};  % 源中心 (969.80,641.89)；下移 3.0 pt 避 AC 虚线（原净距 2.0 pt）
  \node[anchor=center] at (16.758,41.351) {$D'$};  % 源中心 (966.52,670.77)；取 D 与标注撇号合框中心
  % 图注（源页图下居中）
  \node at (34.019,-12.000) {图二};
\end{tikzpicture}
%
% 【自验】
%   1) 探针编译：0 error、0 warning。
%   2) 600 dpi 双渲染差分量测标签净距：B 3.48、A 3.56、D′ 3.64、C 4.07 —— 最小 3.48 pt。
%   3) 线形：AB／AD′／BD′ 实线 0.399 pt；AC／CB／CD′ 虚线节距 on 2.99pt off 2.99pt（源 d=[2.99 2.99]）。
